\documentclass[12pt]{article}
\usepackage{amsmath, amssymb, amsfonts}
\usepackage{geometry}
\usepackage{setspace}
\usepackage{enumitem}
\usepackage{booktabs}

\geometry{margin=1in}
\setstretch{1.2}

\title{ECON 695 -- Problem Set 4}
\date{Due: October 27, 2025}

\begin{document}

\maketitle

\section*{Question 1: Average Treatment Effects}

Suppose we have a random sample of size $N$ of which $N_1 < N$ were randomly assigned to the treatment group $D_i = 1$, and for these individuals we observe $Y_i = Y_i(1)$, while the rest $N_2 = N - N_1$ were assigned to the control group $D_i = 0$, for whom we observe $Y_i = Y_i(0)$.

\begin{enumerate}[label=\arabic*.]

  \item Show that under random assignment of treatment, the ATE ($\alpha = E[Y_i(1) - Y_i(0)]$) is equal to another causal effect of interest, which is Average Treatment on the Treated ($ATT := E[Y_i(1) - Y_i(0) \mid D_i = 1]$). Generally speaking, how would you interpret ATE versus ATT?

    % --- solution to part (1) ---
    \paragraph{Solution (Part 1).}
\textbf{Goal:} Show that, under random assignment of treatment,
\[
\text{ATE} \;:=\; E[Y_i(1)-Y_i(0)] \quad \text{equals} \quad 
\text{ATT} \;:=\; E\!\big[Y_i(1)-Y_i(0)\,\big|\,D_i=1\big].
\]

\textbf{Assumptions.}
\begin{enumerate}[label=(A\arabic*)]
  \item \emph{Random assignment (unconfoundedness):} 
  \((Y_i(1),Y_i(0))\perp D_i\).
  \item \emph{Overlap/positivity:} \(0<\Pr(D_i=1)<1\).
  \item \emph{SUTVA/consistency:} Each unit has well-defined potential outcomes and the observed outcome satisfies
  \(Y_i = D_i\,Y_i(1) + (1-D_i)\,Y_i(0)\).
\end{enumerate}

\noindent\textbf{Key implication of (A1).} For any integrable function \(g\),
\[
E\!\big[g\big(Y_i(1),Y_i(0)\big)\,\big|\,D_i=d\big]
\;=\;
E\!\big[g\big(Y_i(1),Y_i(0)\big)\big],
\qquad d\in\{0,1\}.
\]
In particular,
\[
E\!\big[Y_i(1)\,\big|\,D_i=1\big]=E\!\big[Y_i(1)\big],
\qquad
E\!\big[Y_i(0)\,\big|\,D_i=1\big]=E\!\big[Y_i(0)\big].
\]

\noindent\textbf{Derivation.} Using linearity of expectation,
\[
\text{ATT}
= E\!\big[Y_i(1)-Y_i(0)\,\big|\,D_i=1\big]
= E\!\big[Y_i(1)\,\big|\,D_i=1\big] - E\!\big[Y_i(0)\,\big|\,D_i=1\big].
\]
By random assignment (A1),
\[
E\!\big[Y_i(1)\,\big|\,D_i=1\big] = E\!\big[Y_i(1)\big],
\qquad
E\!\big[Y_i(0)\,\big|\,D_i=1\big] = E\!\big[Y_i(0)\big].
\]
Therefore,
\[
\text{ATT}
= E\!\big[Y_i(1)\big] - E\!\big[Y_i(0)\big]
= E\!\big[Y_i(1)-Y_i(0)\big]
= \text{ATE}.
\]

\noindent\textbf{Interpretation (ATE vs.\ ATT).} 
\begin{itemize}[leftmargin=1.5em]
  \item \(\text{ATE}=E[Y_i(1)-Y_i(0)]\) is the average causal effect in the full population.
  \item \(\text{ATT}=E[Y_i(1)-Y_i(0)\mid D_i=1]\) is the average causal effect among those assigned to treatment.
\end{itemize}
Under random assignment, treated and untreated groups are (in distribution) the same ex ante, so the average effect among the treated equals the population average effect: \(\text{ATT}=\text{ATE}\). Without random assignment, these generally differ because \(D_i\) may be related to \((Y_i(1),Y_i(0))\) (selection).

  \item Show:
  \[
  E[Y_i D_i] = E[Y_i(1) \mid D_i = 1] \times \Pr(D_i = 1)
  \]
  \[
  E[Y_i (1 - D_i)] = E[Y_i(0) \mid D_i = 0] \times \Pr(D_i = 0)
  \]

  % --- solution to part (2) ---

\paragraph{Solution (Part 2).}
We use two ingredients throughout: (i) the Law of Iterated Expectations (LIE),
\(E[W]=E\!\big[\,E[W\mid D_i]\,\big]\), and (ii) consistency/SUTVA,
\(Y_i=D_iY_i(1)+(1-D_i)Y_i(0)\).

\medskip
\noindent\underline{First identity: \(E[Y_i D_i] = E[Y_i(1)\mid D_i=1]\times \Pr(D_i=1)\).}
\[
\begin{aligned}
E[Y_i D_i]
&= E\!\big[\,E[Y_i D_i \mid D_i]\,\big] &&\text{(LIE)} \\
&= E\!\big[\,D_i\,E[Y_i \mid D_i]\,\big] &&\text{(pull out }D_i\text{, measurable w.r.t.\ }D_i\text{)}\\
&= \Pr(D_i=1)\,E[Y_i \mid D_i=1] &&\text{(only the }D_i=1\text{ term survives)}\\
&= \Pr(D_i=1)\,E[Y_i(1) \mid D_i=1] &&\text{(consistency when }D_i=1\text{)}.
\end{aligned}
\]

\medskip
\noindent\underline{Second identity: \(E[Y_i(1-D_i)] = E[Y_i(0)\mid D_i=0]\times \Pr(D_i=0)\).}
\[
\begin{aligned}
E[Y_i(1-D_i)]
&= E\!\big[\,E[Y_i(1-D_i) \mid D_i]\,\big] &&\text{(LIE)} \\
&= E\!\big[(1-D_i)\,E[Y_i \mid D_i]\big] &&\text{(pull out }1-D_i\text{)}\\
&= \Pr(D_i=0)\,E[Y_i \mid D_i=0] &&\text{(only the }D_i=0\text{ term survives)}\\
&= \Pr(D_i=0)\,E[Y_i(0) \mid D_i=0] &&\text{(consistency when }D_i=0\text{)}.
\end{aligned}
\]

\medskip
\noindent These equalities require only LIE and consistency/SUTVA; random assignment is
not needed for part (2).

  \item Show that under random assignment of treatment,
  \[
  \hat{\alpha} = 
  \frac{\sum_{i=1}^N D_i Y_i}{\sum_{i=1}^N D_i}
  - 
  \frac{\sum_{i=1}^N (1 - D_i) Y_i}{\sum_{i=1}^N (1 - D_i)}
  \]
  is a consistent estimator for $\alpha = E[Y_i(1)] - E[Y_i(0)]$.

  % --- solution to part (3) ---

  \paragraph{Solution (Part 3).}
    Write the difference-in-means estimator as
    \[
    \hat{\alpha}
    = \underbrace{\frac{\sum_{i=1}^N D_i Y_i}{\sum_{i=1}^N D_i}}_{\bar Y_1}
    \;-\;
    \underbrace{\frac{\sum_{i=1}^N (1-D_i) Y_i}{\sum_{i=1}^N (1-D_i)}}_{\bar Y_0}
    = \bar Y_1 - \bar Y_0 ,
    \]
    where $\bar Y_1$ and $\bar Y_0$ are the sample means in the treated and control
    groups, respectively.

    \textbf{Assumptions.} (i) Random assignment: $(Y_i(1),Y_i(0))\!\perp\! D_i$; (ii) overlap:
    $0<\Pr(D_i=1)<1$; (iii) SUTVA/consistency: $Y_i = D_iY_i(1) + (1-D_i)Y_i(0)$; (iv) finite first moments.

    \medskip
    \noindent\textbf{Key identities from consistency.}
    \(
    D_i Y_i = D_i Y_i(1), \quad (1-D_i)Y_i = (1-D_i)Y_i(0).
    \)

    \medskip
    \noindent\textbf{Limits of numerators and denominators (LLN).} Let $p:=\Pr(D_i=1)$.
    By the Weak Law of Large Numbers (WLLN),
    \[
    \frac{1}{N}\sum_{i=1}^N D_i \;\xrightarrow{p}\; p,
    \qquad
    \frac{1}{N}\sum_{i=1}^N (1-D_i) \;\xrightarrow{p}\; 1-p.
    \]
    Moreover, using the identities above and random assignment,
    \[
    \frac{1}{N}\sum_{i=1}^N D_i Y_i
    = \frac{1}{N}\sum_{i=1}^N D_i Y_i(1)
    \;\xrightarrow{p}\; E[D_i Y_i(1)]
    = E[D_i]\,E[Y_i(1)]
    = p\,E[Y_i(1)],
    \]
    \[
    \frac{1}{N}\sum_{i=1}^N (1-D_i) Y_i
    = \frac{1}{N}\sum_{i=1}^N (1-D_i) Y_i(0)
    \;\xrightarrow{p}\; E[(1-D_i) Y_i(0)]
    = (1-p)\,E[Y_i(0)].
    \]
    The equalities $E[D_i Y_i(1)] = E[D_i]\,E[Y_i(1)]$ and
    $E[(1-D_i)Y_i(0)] = E[1-D_i]\,E[Y_i(0)]$ follow from $(Y_i(1),Y_i(0))\!\perp\! D_i$.

    \medskip
    \noindent\textbf{Ratios and Slutsky/Continuous Mapping.}
    Hence, using the sample means:
    \begin{align*}
    \bar Y_1
    &= 
    \frac{\tfrac{1}{N}\sum_{i=1}^N D_i Y_i}{\tfrac{1}{N}\sum_{i=1}^N D_i}
    \;\xrightarrow{p}\;
    \frac{p\,E[Y_i(1)]}{p}
    = E[Y_i(1)], \\[1em]
    \bar Y_0
    &=
    \frac{\tfrac{1}{N}\sum_{i=1}^N (1-D_i) Y_i}{\tfrac{1}{N}\sum_{i=1}^N (1-D_i)}
    \;\xrightarrow{p}\;
    \frac{(1-p)\,E[Y_i(0)]}{1-p}
    = E[Y_i(0)].
    \end{align*}

    \noindent Therefore, by Slutsky's theorem,
    \[
    \hat{\alpha}
    = \bar Y_1 - \bar Y_0
    \;\xrightarrow{p}\;
    E[Y_i(1)] - E[Y_i(0)]
    = \alpha.
    \]
    Thus, $\hat{\alpha}$ is a consistent estimator for the ATE under random assignment and overlap.

\end{enumerate}

\newpage

\section*{Question 2: Self Sufficiency Program (see ssp.ipynb)}

\noindent\textbf{Background from Card and Hyslop (2005):}
\begin{itemize}
    \item In the Self Sufficiency Project (SSP), members of a randomly assigned treatment group could receive a subsidy for full-time work. The subsidy was available for 3 years, but only to people who began working full time within 12 months of random assignment.
    \item SSP provides two incentives: (1) find a job and leave welfare within a year, and (2) choose work over welfare in the longer term.
    \item This question is designed to apply the IV-LATE framework to estimate the causal effects of SSP on leaving welfare, and characterize the compliers who would not have found a full-time job in the absence of the SSP subsidy. The data set \texttt{welfare.csv} contains 5,480 observations for people in the SSP experiment. The data set has the following variables:
\end{itemize}

\begin{itemize}
    \item \texttt{treatment} = 1 if in treatment group
    \item \texttt{ft15, ft20, ft24, ft48} = set of dummies if working full time at months 15, 20, 24, and 48 after random assignment
    \item \texttt{welfare15...welfare48} = set of dummies if receiving welfare at months 15, 20, 24, and 48 after random assignment
    \item Other characteristics:
    \begin{itemize}
        \item \texttt{imm} = 1 if immigrant
        \item \texttt{hsgrad} = 1 if high school graduate (or higher) education
        \item \texttt{agelt25} = 1 if age is 25 or less
        \item \texttt{age35p} = 1 if age is 35 or more
        \item \texttt{working\_at\_baseline} = 1 if person was working when assigned to treatment/control group
        \item \texttt{anykidsu6} = 1 if any children under age 6
        \item \texttt{nevermarried} = 1 if never married (all cases are single parents)
    \end{itemize}
\end{itemize}

\begin{enumerate}[label=\arabic*.]

    \item \textbf{(First stage)} Estimate first stage models for the probability of working full time in months 15, 20, 24, and 48, using treatment as the instrumental variable:
    \[
    FT_i(t) = \pi_0 + \pi_1 \text{treatment}_i + \varepsilon_i, \quad \text{for } t = 15, 20, 24, 48
    \]
    What does $\pi_0$ represent? What does $\pi_1$ represent?

    % --- solution to part (1) ---
    \paragraph{Solution (Part 1).}
    The estimated first-stage regression for each time horizon takes the form:
    \[
    FT_i(t) = \pi_0 + \pi_1 \text{treatment}_i + \varepsilon_i, \quad \text{for } t = 15, 20, 24, 48.
    \]

    \textbf{Interpretation:}
    \begin{itemize}
        \item $\pi_0$ represents the \emph{average probability of working full time} among individuals in the control group ($D_i = 0$):
        \[
        \pi_0 = E[FT_i(t) \mid D_i = 0].
        \]
        \item $\pi_1$ captures the \emph{causal effect of treatment assignment on full-time employment}, i.e. the difference in the average probability of working full time between treated and control groups:
        \[
        \pi_1 = E[FT_i(t) \mid D_i = 1] - E[FT_i(t) \mid D_i = 0].
        \]
    \end{itemize}

    \textbf{Empirical Results:}
    The table below summarizes the OLS estimates of $\pi_0$ and $\pi_1$ at each time period $t$:

    \begin{center}
    \begin{tabular}{@{}lcccccc@{}}
    \toprule
    Month ($t$) & $\pi_0$ & $\pi_1$ & SE($\pi_1$) & $t(\pi_1)$ & $N$ \\
    \midrule
    15 & 0.1472 & 0.1374 & 0.011 & 12.55 & 5,480 \\
    20 & 0.1609 & 0.1147 & 0.011 & 10.16 & 5,245 \\
    24 & 0.1602 & 0.1074 & 0.011 & 9.55 & 5,224 \\
    48 & 0.2337 & 0.0506 & 0.013 & 4.01 & 4,796 \\
    \bottomrule
    \end{tabular}
    \end{center}

    \textbf{Discussion:}
    Across all months, the coefficient $\pi_1$ is positive and statistically significant at the 1\% level, indicating that assignment to the SSP treatment group increased the probability of full-time employment relative to the control group. 

    The effect is largest at $t=15$ (about 13.7 percentage points) and declines steadily over time, reaching about 5 percentage points by month 48. This pattern suggests that the SSP program had its strongest impact on promoting full-time work shortly after random assignment, with the effect tapering off in the longer term. 


    \item \textbf{(Reduced form)} Estimate reduced form models for the probability of being on welfare in months 15, 20, 24, and 48, using treatment as the instrumental variable:
    \[
    \text{Welfare}_i(t) = \delta_0 + \delta_1 \text{treatment}_i + \nu_i, \quad \text{for } t = 15, 20, 24, 48
    \]
    where $\text{Welfare}_i(t) = 1$ if person $i$ is on welfare $t$ months since random assignment.

    % --- solution to part (2) ---
    \paragraph{Solution (Part 2).}
    For each horizon \(t\in\{15,20,24,48\}\) we estimate
    \[
    \text{Welfare}_i(t)=\delta_0+\delta_1\,\text{treatment}_i+\nu_i,
    \]
    where \(\text{Welfare}_i(t)=1\) if individual \(i\) is on welfare \(t\) months after random assignment.  
    Here, \(\delta_0=E[\text{Welfare}_i(t)\mid Z_i=0]\) is the control mean (baseline probability of welfare receipt), and \(\delta_1=E[\text{Welfare}_i(t)\mid Z_i=1]-E[\text{Welfare}_i(t)\mid Z_i=0]\) is the \emph{intent-to-treat (ITT)} effect of being offered SSP on welfare status.

    \begin{center}
    \begin{tabular}{@{}lcccccccc@{}}
    \toprule
    $t$ & $\delta_0$ & $\delta_1$ & SE($\delta_1$) & $t(\delta_1)$ & $F(\delta_1)$ & $N$ & $\Pr[W{=}1\mid Z{=}0]$ & $\Pr[W{=}1\mid Z{=}1]$ \\
    \midrule
    15 & 0.810 & $-0.141$ & 0.0117 & $-12.06$ & 145.33 & 5{,}480 & 0.810 & 0.669 \\
    20 & 0.769 & $-0.121$ & 0.0124 & $-9.73$ & 94.73 & 5{,}245 & 0.769 & 0.648 \\
    24 & 0.740 & $-0.108$ & 0.0127 & $-8.49$ & 72.09 & 5{,}224 & 0.740 & 0.632 \\
    48 & 0.592 & $-0.042$ & 0.0143 & $-2.94$ & 8.63 & 4{,}796 & 0.592 & 0.550 \\
    \bottomrule
    \end{tabular}
    \end{center}

    \noindent\emph{Notes.} The last two columns equal the model-implied probabilities: \(\Pr[W{=}1\mid Z{=}0]=\delta_0\) and \(\Pr[W{=}1\mid Z{=}1]=\delta_0+\delta_1\).

    \medskip
    \noindent\textbf{Interpretation.} Assignment to SSP substantially \emph{reduced} welfare receipt at all horizons: \(\delta_1<0\) and statistically significant (1\% level or better). At \(t=15\), the control mean is about \(81\%\) while the treated mean is about \(67\%\), a decline of \(\approx 14\) percentage points (\(t\approx-12.1\)). The ITT effect attenuates over time to \(\approx-12\) p.p.\ at 20 months, \(\approx-11\) p.p.\ at 24 months, and \(\approx-4\) p.p.\ by 48 months (still significant, \(t\approx-2.94\)). This pattern is consistent with the program inducing earlier exits from welfare that partly dissipate once the subsidy period ends.  

    \item \textbf{(2SLS)} Estimate 2SLS models of the equation:
    \[
    \text{Welfare}_i(t) = \beta_0 + \beta_1 FT_i(t) + u_i, \quad \text{for } t = 15, 20, 24, 48
    \]
    using treatment as an instrument for $FT_i(t)$. Verify that the 2SLS estimate in each case is the ratio of the reduced form and first stage coefficients. \emph{Hint:} use \texttt{linearmodels.IV2SLS}.

    % --- solution to part (3) ---
    \paragraph{Solution (Part 3).}
    We estimate, for each horizon \(t\in\{15,20,24,48\}\),
    \[
    \text{Welfare}_i(t) \;=\; \beta_0 \;+\; \beta_1\, FT_i(t) \;+\; u_i,
    \]
    using \(\text{treatment}_i\) as an instrument for the endogenous \(FT_i(t)\).
    Identification relies on: (i) random assignment (independence); (ii) instrument
    relevance (first-stage \(\pi_1\neq 0\)); (iii) exclusion (assignment affects
    welfare only through \(FT\)); and (iv) monotonicity (no defiers). With a binary
    instrument and treatment, \(\beta_1\) is the LATE/CACE: the average causal effect
    of full-time work on welfare among compliers at horizon \(t\).

    \medskip
    \noindent\textbf{2SLS estimates.} (HC1 standard errors)

    \begin{center}
    \begin{tabular}{@{}lccccc@{}}
    \toprule
    Month ($t$) & $\beta_0$ & $\beta_1$ & SE($\beta_1$) & $t(\beta_1)$ & $N$ \\
    \midrule
    15 & 0.960872 & $-1.026628$ & 0.082099 & $-12.504746$ & 5{,}480 \\
    20 & 0.939023 & $-1.055033$ & 0.105477 & $-10.002532$ & 5{,}245 \\
    24 & 0.901966 & $-1.007911$ & 0.116582 & $-8.645537$  & 5{,}224 \\
    48 & 0.785521 & $-0.829871$ & 0.259718 & $-3.195279$  & 4{,}796 \\
    \bottomrule
    \end{tabular}
    \end{center}

    \noindent\textbf{Verification.} For each \(t\), we verified in the notebook that
    \[
    \widehat{\beta}_1 \;=\; \frac{\widehat{\delta}_1}{\widehat{\pi}_1},
    \]
    i.e., the 2SLS slope equals the ratio of the reduced-form coefficient on
    \(\text{treatment}\) in the welfare equation (\(\widehat{\delta}_1\)) to the
    first-stage coefficient of \(\text{treatment}\) in the \(FT\) equation
    (\(\widehat{\pi}_1\)); the numerical equality holds up to rounding.

    \medskip
    \noindent\textbf{Interpretation.}
    The estimates imply a large negative LATE of full-time work on welfare receipt
    for compliers:
    at \(t=15\)–\(24\) months, \(\widehat{\beta}_1\approx -1\) and highly significant,
    suggesting that moving from not working FT to working FT (for those whose FT
    status is shifted by the SSP offer) nearly eliminates welfare participation at
    those horizons. By \(t=48\), the effect attenuates to about \(-0.83\) and remains
    statistically significant at conventional levels. This time pattern mirrors the
    program’s temporary subsidy window and the weakening first stage at \(t=48\).
    (Note: with a linear probability model, sampling noise can push \(\widehat{\beta}_1\)
    slightly outside \([-1,1]\); the substantive conclusion is unchanged.)


    \item \textbf{(Characteristics of Compliers)} 
        Find the mean characteristics of the compliers in month $t = 15$: i.e., means of the variables 
        \texttt{imm}, \texttt{hsgrad}, \texttt{agelt25}, \texttt{age35p}, 
        \texttt{working\_at\_baseline}, \texttt{anykidsu6}, and \texttt{nevermarried}. 
        Compare these to the characteristics of always takers and never takers in month 15. 
        \emph{Hint:} define outcome 
        \[
        Y_i = X_i \times FT_i(15).
        \]

    % --- solution to part (4) ---
    \paragraph{Solution (Part 4).}
    For \(t = 15\), we compute the mean characteristics of the \emph{compliers} and compare them to those of \emph{always takers} and \emph{never takers}. Following the IV--LATE framework, let \(Z_i\) denote assignment to treatment and \(FT_i(15)\) the indicator for working full time at 15 months. The groups are defined as:
    \[
    \text{Always takers: } E[X_i \mid FT_i(15) = 1,\, Z_i = 0],
    \quad
    \text{Never takers: } E[X_i \mid FT_i(15) = 0,\, Z_i = 1],
    \]
    \[
    \text{Compliers: }
    E[X_i \mid \text{complier}] 
    = 
    \frac{E[X_i \mid Z_i=1] - E[X_i \mid Z_i=0]}
    {P(FT_i(15)=1 \mid Z_i=1) - P(FT_i(15)=1 \mid Z_i=0)}.
    \]

    Using the sample shares
    \[
    P(FT_i(15)=1\mid Z_i=0)=0.148, \quad 
    P(FT_i(15)=1\mid Z_i=1)=0.285,
    \]
    we obtain the following estimates:

    \begin{center}
    \begin{tabular}{@{}lccc@{}}
    \toprule
    Characteristic & Compliers & Always Takers & Never Takers \\
    \midrule
    imm & $-0.026$ & $0.078$ & $0.139$ \\
    hsgrad & $0.117$ & $0.615$ & $0.414$ \\
    agelt25 & $-0.058$ & $0.172$ & $0.170$ \\
    age35p & $0.033$ & $0.298$ & $0.327$ \\
    working\_at\_baseline & $-0.048$ & $0.468$ & $0.122$ \\
    anykidsu6 & $-0.145$ & $0.502$ & $0.520$ \\
    nevermarried & $-0.011$ & $0.485$ & $0.483$ \\
    \bottomrule
    \end{tabular}
    \end{center}

    \noindent\textbf{Interpretation.}
    \begin{itemize}[leftmargin=1.5em]
        \item \textbf{Compliers} lie between always and never takers, consistent with the IV--LATE interpretation that the SSP affects a marginal subgroup. They are less likely to be immigrants or to have young children, and have lower baseline employment compared to always takers. Their education (\texttt{hsgrad}) is lower than that of always takers but higher than never takers.
        \item \textbf{Always takers} are the most advantaged group—older (\texttt{age35p}), better educated, and more likely to be working at baseline—implying that they would work full time regardless of treatment.
        \item \textbf{Never takers} are the most constrained—less educated, lower baseline work rates, and more likely to have young children—explaining why they do not respond to the SSP offer.
    \end{itemize}

    \noindent Overall, the SSP subsidy appears to induce full-time employment primarily among \emph{moderately disadvantaged individuals}—those not as attached to the labor market as always takers, but facing fewer barriers than never takers. This pattern aligns with the local average treatment effect (LATE) interpretation, identifying the program’s causal impact among compliers.

\end{enumerate}



\end{document}
